\subsection{Acquisition et formation d'images SAR}

Les systèmes de radar à synthèse d'ouverture (SAR) reposent sur l'acquisition cohérente d'échos radar le long d'une trajectoire synthétique, permettant de reconstruire une image à haute résolution indépendamment de la distance de la cible. La phase d'acquisition conditionne la qualité radiométrique, la résolution et le volume de données à traiter.

\subsubsection{Principes fondamentaux de l’acquisition SAR}

Les travaux fondateurs de Bamler et Hartl \cite{Bamler1993_SAR} formalisent le lien entre la géométrie de vol, la modulation du signal émis et la phase mesurée, posant les bases des schémas de focalisation. La cohérence de phase entre impulsions successives, garantie par une horloge stable et une synchronisation fine, est essentielle à la synthèse d’ouverture.

\subsubsection{Architectures d’acquisition}

Plusieurs architectures se distinguent :
\begin{itemize}
  \item \textbf{SAR à impulsions (Pulse SAR)} : architecture des grands systèmes spatiaux/aéroportés ; forte portée et dynamique, au prix d’une électronique d’émission puissante et d’une chaîne RF complexe \cite{Krieger2010_InterferometricSAR}.
  \item \textbf{SAR FMCW} : adapté aux plateformes légères (UAV) ; l’onde chirpée est déchirpée en réception, réduisant la bande utile à numériser. Des démonstrateurs complets et légers sont documentés dans \cite{Aguasca2013_ARBRES_UAV_SAR,Svedin2021_UAVSAR}.
  \item \textbf{SAR passif} : le capteur n’émet pas, exploite des illuminateurs d’opportunité ; consommation minimale mais corrélation et synchronisation plus exigeantes \cite{Sun2020_PassiveUAVSAR}.
  \item \textbf{Multi-canaux / polarimétrique} : enrichissement informationnel (HH/HV/VV) au prix d’une stricte synchronisation inter-voies et d’un volume de données accru \cite{Krieger2010_InterferometricSAR}.
\end{itemize}

\subsubsection{Chaîne d’acquisition embarquée}

La chaîne moderne numérisée intègre : (i) un synthétiseur (DDS/PLL) pour générer le chirp, (ii) un émetteur à puissance modulable, (iii) une réception LNA--mélangeur--IF--ADC, et (iv) une électronique de contrôle (FPGA/SBC) pour tramage et stockage. Sur UAV FMCW, des systèmes de quelques kilogrammes réalisent des acquisitions cohérentes avec navigation intégrée (GNSS/IMU) \cite{Aguasca2013_ARBRES_UAV_SAR,Svedin2021_UAVSAR}. Les formats de métadonnées (SigMF) facilitent l’archivage et le post-traitement \cite{Balister2017_SigMF}.

\subsubsection{Tendances actuelles}

Les approches \textit{Software Defined} et les SoC RF intégrés permettent reconfiguration en vol, compacité et faible latence ; elles s’appuient sur des horloges distribuées (10~MHz/1~PPS) et des liens déterministes (JESD204B/C) \cite{AMD_RFSoC_DS889,ADI_JESD204B_Survival}. L’optimisation énergétique et la génération d’aperçus (range--Doppler) à bord sont très étudiées sur UAV \cite{Svedin2021_UAVSAR,Sun2020_PassiveUAVSAR}.

\vspace{0.5em}
\noindent\textit{En résumé}~: l’acquisition SAR a évolué d’architectures impulsionnelles lourdes vers des systèmes FMCW/SDR miniaturisés, cohérents et reconfigurables, adaptés aux contraintes de masse et d’énergie des UAV \cite{Aguasca2013_ARBRES_UAV_SAR,Svedin2021_UAVSAR}.

\subsection{Cartes d'acquisition pour SAR embarqué : état de l'art}

L’acquisition SAR exige des convertisseurs rapides et cohérents en phase, une chaîne d’horloges stable (10~MHz~+~1~PPS) et une voie d’horodatage fiable. Quatre familles dominent la littérature et l’offre industrielle : (1) \textit{SDR} reconfigurables, (2) SoC RF intégrant ADC/DAC (\textit{RFSoC}), (3) numériseurs PCIe/PXIe, et (4) mezzanines \textit{FMC/FMC+} sur FPGA.

\paragraph{(1) SDR reconfigurables.}
Les USRP X310 offrent bande de base élevée (jusqu’à \textasciitilde160~MHz selon carte RF), FPGA et cohérence via 10~MHz/1~PPS (GPSDO) \cite{Ettus_X310_Spec}. Adapté aux SAR FMCW et aux démonstrateurs aéroportés souples.

\paragraph{(2) SoC RF intégrés (\textit{RFSoC}).}
Les Zynq UltraScale+ RFSoC intègrent ADC/DAC multi-voies et logique programmable, combinant compacité, faible latence et synchronisation fine pour architectures tout-en-un \cite{AMD_RFSoC_DS889}.

\paragraph{(3) Numériseurs PCIe/PXIe.}
Des cartes comme NI PXIe-5775 (FlexRIO) apportent des GSa/s avec horloges de châssis, utiles pour bancs et pods haute bande instantanée \cite{NI_PXIe5775_Spec}.

\paragraph{(4) Mezzanines \textit{FMC/FMC+}.}
Des modules comme Abaco FMC134 (JESD204B/C) permettent un choix fin des convertisseurs et du FPGA porteur pour des prototypes performants et compacts \cite{Abaco_FMC134}. La distribution d’horloge et la latence déterministe JESD204B/C sont clés \cite{ADI_JESD204B_Survival}.

\subsubsection*{Critères de sélection}
\begin{itemize}
  \item \textbf{Bande instantanée et $f_s$} : FMCW déchirpé (quelques--dizaines de MSa/s) vs impulsionnel large bande (centaines de MSa/s à GSa/s).
  \item \textbf{ENOB, SNR, SFDR} : dynamique et qualité de focalisation.
  \item \textbf{Synchronisation} : 10~MHz, 1~PPS, MCS/JESD204 déterministe, \textit{phase noise} faible.
  \item \textbf{Chaîne de données} : buffers, DMA, lien hôte/SSD (PCIe/10~GigE/NVMe).
  \item \textbf{Contraintes UAV} : masse, consommation, thermique/vibrations.
\end{itemize}

\subsection{Cartes d'acquisition pour SAR embarqué : état de l'art}

L’acquisition SAR exige des convertisseurs rapides et cohérents en phase, une chaîne d’horloges stable (10~MHz~+~1~PPS) et une voie d’horodatage fiable. Quatre familles dominent la littérature et l’offre industrielle : (1) \textit{SDR} reconfigurables, (2) SoC RF intégrant ADC/DAC (\textit{RFSoC}), (3) numériseurs PCIe/PXIe, et (4) mezzanines \textit{FMC/FMC+} sur FPGA.

\paragraph{(1) SDR reconfigurables.}
Les USRP X310 offrent bande de base élevée (jusqu’à \textasciitilde160~MHz selon carte RF), FPGA et cohérence via 10~MHz/1~PPS (GPSDO) \cite{Ettus_X310_Spec}. Adapté aux SAR FMCW et aux démonstrateurs aéroportés souples.

\paragraph{(2) SoC RF intégrés (\textit{RFSoC}).}
Les Zynq UltraScale+ RFSoC intègrent ADC/DAC multi-voies et logique programmable, combinant compacité, faible latence et synchronisation fine pour architectures tout-en-un \cite{AMD_RFSoC_DS889}.

\paragraph{(3) Numériseurs PCIe/PXIe.}
Des cartes comme NI PXIe-5775 (FlexRIO) apportent des GSa/s avec horloges de châssis, utiles pour bancs et pods haute bande instantanée \cite{NI_PXIe5775_Spec}.

\paragraph{(4) Mezzanines \textit{FMC/FMC+}.}
Des modules comme Abaco FMC134 (JESD204B/C) permettent un choix fin des convertisseurs et du FPGA porteur pour des prototypes performants et compacts \cite{Abaco_FMC134}. La distribution d’horloge et la latence déterministe JESD204B/C sont clés \cite{ADI_JESD204B_Survival}.

\subsubsection*{Critères de sélection}
\begin{itemize}
  \item \textbf{Bande instantanée et $f_s$} : FMCW déchirpé (quelques--dizaines de MSa/s) vs impulsionnel large bande (centaines de MSa/s à GSa/s).
  \item \textbf{ENOB, SNR, SFDR} : dynamique et qualité de focalisation.
  \item \textbf{Synchrox    nisation} : 10~MHz, 1~PPS, MCS/JESD204 déterministe, \textit{phase noise} faible.
  \item \textbf{Chaîne de données} : buffers, DMA, lien hôte/SSD (PCIe/10~GigE/NVMe).
  \item \textbf{Contraintes UAV} : masse, consommation, thermique/vibrations.
\end{itemize}
\paragraph*{Transition vers le stockage des données.}
Ainsi, le choix de la carte d’acquisition détermine non seulement les performances de la chaîne analogique et numérique, mais aussi le \textit{débit brut} généré par le système. La fréquence d’échantillonnage, le nombre de canaux et la résolution effective fixent directement le volume d’information à enregistrer. Ces paramètres conditionnent donc les besoins en mémoire tampon, en capacité de stockage et en gestion d’énergie, qui constituent les défis majeurs de la section suivante consacrée au \textbf{stockage des données SAR embarquées}.